% CLI Reference: Subcommands, Options, TOML Configuration, and Common Workflows
% Source: docs/source/cli.md (migrated to LaTeX)
% Uses macros from preamble.tex

\section{CLI Reference}
\label{sec:cli-reference}

TIDAL provides the \texttt{tidal} command-line tool with nine subcommands for the full derive-to-measure workflow.
\textbf{Zero additional dependencies} --- uses only \texttt{argparse} and \texttt{tomllib} from the Python standard library.

\subsection{Installation}

\begin{lstlisting}[style=bash]
# Install with uv (recommended)
uv sync --all-extras

# Or via pip
pip install -e ".[dev]"

# Verify
tidal --version
\end{lstlisting}

\subsection{Subcommands}

\subsubsection{\texttt{tidal derive} --- Derive equations from a Lagrangian}

Generates a Wolfram Language script (\texttt{.wls}) from a TOML configuration file and optionally runs it with \texttt{wolframscript}.

\begin{lstlisting}[style=bash]
# Generate .wls only
tidal derive examples/scalar_field/theory.toml

# Generate .wls and run wolframscript via the generated run.sh
bash examples/scalar_field/run.sh
\end{lstlisting}

\subsubsection{\texttt{tidal simulate} --- Run a PDE simulation}

Loads a JSON equation specification and runs a numerical simulation with configurable parameters.

\begin{lstlisting}[style=bash]
# Basic simulation with defaults
tidal simulate examples/data/klein_gordon_1d.json

# With parameter overrides and custom settings
tidal simulate examples/data/coupled_scalars.json \
    --param m_phi2=1.0 --param g=0.5 \
    --t-end 20.0 --dt 0.005 \
    --grid-shape 256 \
    --ic gaussian --ic-width 5.0 \
    --bc periodic \
    --scheme scipy

# Constraint-only solve (no time evolution)
tidal simulate examples/data/chern_simons.json --mode constraint
\end{lstlisting}

\paragraph{Options}

\begin{table}[htbp]
\centering
\begin{tabular}{@{}llp{4.5cm}@{}}
\toprule
Flag & Default & Description \\
\midrule
\texttt{--param KEY=VALUE}          & (from JSON)       & Override symbolic coefficient \\
\texttt{--t-end FLOAT}             & 10.0              & Simulation end time \\
\texttt{--dt FLOAT}                & 0.01              & Time step \\
\texttt{--grid-shape INT}          & 128               & Grid points per axis \\
\texttt{--ic \{gaussian,...\}}      & gaussian          & Initial condition preset \\
\texttt{--ic-width FLOAT}          & 5.0               & Gaussian pulse width \\
\texttt{--ic-component FIELD}      & (first field)     & Field component for IC (multi-field systems) \\
\texttt{--ic-center X[,Y,Z]}      & (domain center)   & Gaussian center position \\
\texttt{--ic-amplitude FLOAT}      & 1.0               & Gaussian amplitude \\
\texttt{--ic-wavevector K[,K,K]}   & ---               & Wavevector for travelling wave packets \\
\texttt{--ic-formula EXPR}         & ---               & Custom IC formula (Python expression with \texttt{x}, \texttt{y}, \texttt{z}, \texttt{np}) \\
\texttt{--ic-formula-velocity EXPR}& ---               & Velocity formula for \texttt{--ic formula} (enables travelling waves) \\
\texttt{--ic-field FIELD:EXPR}     & ---               & Per-field IC override formula (repeatable; \texttt{FIELD:velocity:EXPR} for velocity) \\
\texttt{--ic-file PATH}            & ---               & Load IC from \texttt{.npy} file or simulation output directory \\
\texttt{--ic-noise-seed N}         & ---               & Random seed for reproducible \texttt{--ic noise} \\
\texttt{--bc TYPES}                & periodic          & Boundary conditions (comma-separated per axis) \\
\texttt{--scheme \{ida,...,auto\}}  & auto              & Solver scheme \\
\texttt{--mode \{evolve,constraint\}}& evolve           & Simulation mode \\
\texttt{--plot} / \texttt{--no-plot}& \texttt{--plot}   & Enable/disable plotting \\
\texttt{--output PATH}             & ---               & Output path (directory for snapshot data; image extension for plot-only) \\
\bottomrule
\end{tabular}
\end{table}

\subsubsection{\texttt{tidal inspect} --- Display equation system info}

Shows a summary of a JSON equation specification without running a simulation.

\begin{lstlisting}[style=bash]
tidal inspect examples/data/chern_simons.json
\end{lstlisting}

Output includes: spacetime dimension, field names, equation count, operator types, mass/coupling matrices.

\subsubsection{\texttt{tidal list} --- Discover available JSON specs}

Searches for JSON specification files in \texttt{examples/data/}.

\begin{lstlisting}[style=bash]
tidal list
\end{lstlisting}

\subsubsection{\texttt{tidal validate} --- Validate a JSON specification}

Checks that a JSON file conforms to the expected equation system schema.

\begin{lstlisting}[style=bash]
tidal validate examples/data/klein_gordon_1d.json

# JSON output for scripting
tidal validate examples/data/klein_gordon_1d.json --json
\end{lstlisting}

\subsubsection{\texttt{tidal measure} --- Extract physics measurements}

Loads a snapshot directory produced by \texttt{tidal simulate --output <dir>} and runs measurement analyses: energy decomposition, conversion probability, mixing length, spectral analysis, and conservation diagnostics.

\begin{lstlisting}[style=bash]
# Full summary (energy + conservation + auto-detect conversion + mixing)
tidal measure result_dir/ --spec spec.json

# Specific measurements
tidal measure result_dir/ --what conversion,mixing \
    --source phi_0 --target chi_0

# JSON output for scripting
tidal measure result_dir/ --what energy,conservation --json

# Save 2x3 measurement plot
tidal measure result_dir/ --output measurements.png
\end{lstlisting}

The JSON spec can be auto-discovered from \texttt{metadata.json} in the snapshot directory (stored by \texttt{tidal simulate}) or provided explicitly via \texttt{--spec}.

\paragraph{Options}

\begin{table}[htbp]
\centering
\begin{tabular}{@{}llp{5cm}@{}}
\toprule
Flag & Default & Description \\
\midrule
\texttt{--spec PATH}            & (auto-discovered) & JSON equation specification \\
\texttt{--what TYPES}           & summary           & Measurement types (comma-separated) \\
\texttt{--source FIELDS}        & (auto-detect)     & Source field(s) for conversion (comma-separated) \\
\texttt{--target FIELDS}        & (auto-detect)     & Target field(s) for conversion (comma-separated) \\
\texttt{--param KEY=VALUE}      & (from metadata)   & Override parameter from \texttt{metadata.json} \\
\texttt{--energy-threshold T}   & 1e-3              & Conservation threshold \\
\texttt{--output PATH}          & ---               & Save plot (\texttt{.png}/\texttt{.pdf}) \\
\texttt{--json}                 & ---               & JSON output instead of text \\
\texttt{--quiet}                & ---               & Suppress progress messages \\
\bottomrule
\end{tabular}
\end{table}

\textbf{Available measurement types:} \texttt{summary}, \texttt{energy}, \texttt{conversion}, \texttt{mixing}, \texttt{spectrum}, \texttt{spectral\_conversion}, \texttt{dispersion}, \texttt{conservation}, \texttt{effective\_mass}, \texttt{asymptotic}, \texttt{peak\_conversion}, \texttt{velocity}, \texttt{resonance}.

\subsubsection{\texttt{tidal sweep} --- Run parameter sweeps and convergence studies}

Automates running \texttt{simulate} + \texttt{measure} across a parameter grid, collecting scalar metrics into a portable CSV/JSON results table.

\begin{lstlisting}[style=bash]
# 1D sweep: coupling strength
tidal sweep examples/data/coupled_scattering.json \
  --sweep "g0=0.1:0.9:10" \
  --measure conversion,mixing --source phi_0 --target chi_0 \
  --grid-shape 256 --bounds=-40:40 --periodic \
  --ic gaussian --ic-component phi_0 --ic-center=-20 --ic-width 3 \
  --ic-wavevector 3 \
  --t-end 60 --output sweeps/coupling_scan

# 2D sweep: coupling x mass (cartesian product)
tidal sweep examples/data/coupled_scattering.json \
  --sweep "g0=0.1:0.9:5" --sweep "mChi2=0.5:4.0:5" \
  --measure conversion --source phi_0 --target chi_0 \
  --output sweeps/mass_coupling_2d

# Grid convergence study
tidal sweep examples/data/coupled_scattering.json \
  --converge "32,64,128,256,512" \
  --measure conservation \
  --output sweeps/convergence
\end{lstlisting}

\paragraph{Sweep specification syntax}

\begin{table}[htbp]
\centering
\begin{tabular}{@{}lll@{}}
\toprule
Format & Description & Example \\
\midrule
\texttt{PARAM=START:STOP:N}       & $N$ linearly-spaced values & \texttt{g0=0.1:1.0:10} \\
\texttt{PARAM=START:STOP:N:log}   & $N$ log-spaced values      & \texttt{m2=0.01:100:20:log} \\
\texttt{PARAM=V1,V2,...}          & Explicit values            & \texttt{g0=0.1,0.5,0.9} \\
\bottomrule
\end{tabular}
\end{table}

Multiple \texttt{--sweep} flags produce a cartesian product grid.

\paragraph{Output structure}

\begin{lstlisting}
sweeps/coupling_scan/
|- sweep.json          # Sweep definition + provenance
|- results.csv         # Self-contained metrics table
|- results.json        # Same data in JSON
|- g0_0.100/           # Individual run output
|   |- metadata.json
|   |- times.npy
|   |- phi_0.npy, v_phi_0.npy, ...
|- g0_0.200/
|   +-- ...
\end{lstlisting}

The CSV is \textbf{self-contained and portable}: every row includes all parameters (swept + fixed), simulation settings, and measured metrics so the data can be analyzed outside TIDAL.

\paragraph{Options}

\begin{table}[htbp]
\centering
\begin{tabular}{@{}llp{5cm}@{}}
\toprule
Flag & Default & Description \\
\midrule
\texttt{--sweep SPEC}             & ---  & Parameter sweep specification (repeatable) \\
\texttt{--converge SIZES}         & ---  & Grid sizes for convergence study \\
\texttt{--config PATH}            & ---  & TOML sweep configuration file \\
\texttt{--measure TYPES}          & summary & Measurement types (comma-separated) \\
\texttt{--source FIELDS}          & ---  & Source field(s) for conversion measurements \\
\texttt{--target FIELDS}          & ---  & Target field(s) for conversion measurements \\
\texttt{--resume}                 & ---  & Skip completed runs (checks for \texttt{metadata.json}) \\
\texttt{--parallel N}             & 1    & Number of parallel workers \\
\texttt{--adaptive-metric M}      & ---  & Metric driving adaptive refinement \\
\texttt{--adaptive-budget N}      & 20   & Max total points for adaptive sampling \\
\texttt{--adaptive-initial N}     & 5    & Initial coarse grid points \\
\texttt{--adaptive-threshold T}   & 0.01 & Stop when max interval score $< T$ \\
\texttt{--dry-run}                & ---  & Print sweep plan without executing \\
\texttt{--force-large-sweep}      & ---  & Allow sweeps with $> 10{,}000$ runs \\
\bottomrule
\end{tabular}
\end{table}

All \texttt{tidal simulate} flags (\texttt{--param}, \texttt{--grid-shape}, \texttt{--bounds}, \texttt{--ic}, etc.)\ are passed through unchanged.

\paragraph{TOML configuration}

Use \texttt{--config} for reproducible, version-controlled sweep definitions:

\begin{lstlisting}[style=bash]
tidal sweep --config examples/coupled_scattering/sweep_coupling.toml
\end{lstlisting}

CLI flags override TOML values, so \texttt{--config sweep.toml --param g0=0.3} uses TOML for everything except \texttt{g0}. See \texttt{examples/coupled\_scattering/sweep\_advanced.toml} for the full TOML schema.

\paragraph{Adaptive sampling}

Automatically refines near sharp features:

\begin{lstlisting}[style=bash]
tidal sweep spec.json --sweep "g0=0.01:0.9:5" \
  --adaptive-metric P_max --adaptive-budget 20 \
  --adaptive-threshold 0.01 \
  --measure conversion --source phi_0 --target chi_0
\end{lstlisting}

Or via TOML:

\begin{lstlisting}[style=toml]
[sweep.g0]
start = 0.01
stop = 0.9
scale = "adaptive"
initial_count = 5
max_count = 20
metric = "P_max"
threshold = 0.01
\end{lstlisting}

\paragraph{Visualizing sweep results}

Use \texttt{tidal plot}:

\begin{lstlisting}[style=bash]
# 1D sweep -> line plot
tidal plot sweeps/coupling_scan/ --type sweep --metric P_max

# 2D sweep -> heatmap
tidal plot sweeps/mass_coupling_2d/ --type sweep --metric P_max

# Multi-metric subplots
tidal plot sweeps/coupling_scan/ --type sweep --metric P_max,L_mix

# Overlay timeseries from each run
tidal plot sweeps/coupling_scan/ --type sweep-compare --metric conversion

# Convergence -> log-log with fitted order
tidal plot sweeps/convergence/ --type convergence \
    --metric max_energy_error

# Parallel coordinates (multi-parameter exploration)
tidal plot sweeps/output/ --type sweep-parallel --metric P_max

# Tornado chart (parameter importance ranking)
tidal plot sweeps/output/ --type sweep-tornado --metric P_max

# Scatter matrix (pairwise parameter relationships)
tidal plot sweeps/output/ --type sweep-scatter --metric P_max
\end{lstlisting}

\subsubsection{\texttt{tidal analyze} --- Sensitivity analysis on sweep results}

Post-hoc sensitivity analysis on completed parameter sweeps. Determines which parameters most influence a given metric.

\begin{lstlisting}[style=bash]
# Sobol global sensitivity indices (requires SALib package)
tidal analyze sweeps/output/ --sensitivity sobol --metric P_max \
    --bootstrap 100

# Morris screening (simpler, no extra dependencies)
tidal analyze sweeps/output/ --sensitivity morris --metric P_max
\end{lstlisting}

\paragraph{Options}

\begin{table}[htbp]
\centering
\begin{tabular}{@{}llp{5cm}@{}}
\toprule
Flag & Default & Description \\
\midrule
\texttt{--sensitivity METHOD} & ---           & Analysis method: \texttt{sobol} or \texttt{morris} \\
\texttt{--metric METRIC}      & (auto-detect) & Which metric to analyze \\
\texttt{--bootstrap N}        & 100           & Bootstrap samples for Sobol confidence intervals \\
\bottomrule
\end{tabular}
\end{table}

\subsection{TOML Configuration (\texttt{theory.toml})}

The \texttt{tidal derive} command reads a TOML file that defines the field theory:

\begin{lstlisting}[style=toml]
[spacetime]
dimension = 3          # 2 = 1+1D, 3 = 2+1D, 4 = 3+1D
signature = [-1, 1, 1]
coordinates = ["t", "x", "y"]

[fields]
names = ["phi"]
ranks = [0]            # 0 = scalar, 1 = vector, 2 = tensor

[constants]
names = ["m2"]
values = [1.0]

[lagrangian]
expression = "1/2 CD[-a][phi[]] eta[a,b] CD[-b][phi[]] - 1/2 m2 phi[]^2"

[parameters]
output_json = "examples/data/my_theory.json"

# Optional: diagonal metric for curvilinear coordinates
[metric]
diagonal = [-1, 1, "x[]^2"]
\end{lstlisting}

\subsubsection{Linearization (\texttt{[linearization]})}

For theories defined by linearizing a tensor expression (e.g., linearized gravity via xPert), use \texttt{[linearization]} \textbf{instead of} \texttt{[lagrangian]}. The two sections are mutually exclusive.

\begin{lstlisting}[style=toml]
[linearization]
# Tensor expression to linearize (uses CD placeholder, auto-prefixed)
expression = "Einstein[CD][-a, -b]"
# Which [[fields]] entry is the metric perturbation
perturbation_field = "h"
\end{lstlisting}

This generates xPert setup (\texttt{SetupMetricPerturbation}), linearization (\texttt{LinearizeTensorExpression}), and automatic notation conversion before passing to the standard decomposition pipeline.

\paragraph{Example} (\texttt{examples/gravitational\_waves/theory.toml}):

\begin{lstlisting}[style=toml]
[theory]
name = "Linearized Gravity"

[spacetime]
dimension = 4
metric = "minkowski"

[[fields]]
name = "h"
type = "tensor"
rank = 2
symmetry = "symmetric"

[linearization]
expression = "Einstein[CD][-a, -b]"
perturbation_field = "h"

[output]
path = "../data/linearized_gravity.json"
\end{lstlisting}

\subsubsection{Derived Fields}

For theories that need intermediate tensor definitions (e.g., the electromagnetic field strength tensor), use \texttt{[[derived\_fields]]}:

\begin{lstlisting}[style=toml]
[[derived_fields]]
name = "F"
rank = 2
symmetry = "antisymmetric"
definition = "CD[-a][A[-b]] - CD[-b][A[-a]]"
\end{lstlisting}

This generates \texttt{DefTensor} + \texttt{MakeRule} in the \texttt{.wls} output, with automatic substitution before decomposition.

\subsection{Common Workflows}

\subsubsection{Derive and simulate a scalar field theory}

\begin{lstlisting}[style=bash]
# 1. Write theory.toml (see examples/scalar_field/theory.toml)

# 2. Derive equations
tidal derive examples/scalar_field/theory.toml

# 3. Inspect the result
tidal inspect examples/data/klein_gordon_1d.json

# 4. Simulate
tidal simulate examples/data/klein_gordon_1d.json --t-end 20 --ic gaussian
\end{lstlisting}

\subsubsection{Parameter sweep}

\begin{lstlisting}[style=bash]
# Automated sweep (recommended)
tidal sweep examples/data/coupled_scalars.json \
    --sweep "mPhi2=0.5:4.0:8" \
    --measure conversion --source phi_0 --target chi_0 \
    --output sweeps/mass_scan

# Or manually with a shell loop
for m2 in 0.5 1.0 2.0 4.0; do
    tidal simulate examples/data/coupled_scalars.json \
        --param mPhi2=$m2 --output "sweep_m2_${m2}.png"
done
\end{lstlisting}

\subsubsection{Solver schemes}

\begin{itemize}
  \item \texttt{--scheme auto} (default): auto-selects best solver based on equation structure
  \item \texttt{--scheme cvode}: SUNDIALS/CVODE --- adaptive BDF/Adams for wave systems
  \item \texttt{--scheme ida}: SUNDIALS/IDA --- handles all equation types including algebraic constraints
  \item \texttt{--scheme leapfrog}: symplectic St\"ormer--Verlet --- fixed $\Delta t$, zero energy drift
  \item \texttt{--scheme scipy}: \texttt{scipy.integrate.solve\_ivp} --- DOP853/RK45/Radau/BDF
\end{itemize}

\subsubsection{Boundary conditions}

Specify per-axis, comma-separated:

\begin{lstlisting}[style=bash]
# Periodic in x, Neumann in y
tidal simulate spec.json --bc periodic,neumann
\end{lstlisting}

Available types: \texttt{periodic}, \texttt{neumann}, \texttt{dirichlet}.

\subsubsection{Full derive-to-measure pipeline}

\begin{lstlisting}[style=bash]
# 1. Derive equations from a Lagrangian
tidal derive examples/coupled_scalars/theory.toml

# 2. Simulate and save to snapshot directory
tidal simulate examples/data/coupled_scalars.json \
    --param mPhi2=1.0 --param mChi2=4.0 --param gCpl=0.5 \
    --t-end 20.0 --output coupled_scalars_output

# 3. Extract measurements: conversion probability and mixing length
tidal measure coupled_scalars_output/ \
    --what conversion,mixing \
    --source phi_0 --target chi_0

# 4. Save measurement plot
tidal measure coupled_scalars_output/ --output measurements.png
\end{lstlisting}
